\documentclass[12pt]{article}

% Packages
\usepackage[a4paper, margin=1in]{geometry} % Adjust margins
\usepackage{setspace} % For line spacing
\usepackage{tocloft} % Optional: better spacing for TOC
\usepackage[backend=biber,style=apa]{biblatex} % Load biblatex with APA style
\addbibresource{references.bib} % Link to your .bib file

% Set line spacing
\onehalfspacing % or use \doublespacing for double spacing

% Title information
\title{\textbf{\large{Preferences on Tax Enforcement and Tax Administration}}}
\author{\normalsize{Tim Kircali, Stefanie Stantcheva, Matthias Weber}}
\date{\normalsize{30.04.2025}}
\begin{document}
\maketitle
\pagenumbering{arabic}
\section*{Related Literature}
First, this study contributes to the growing body of research that measures support for economic policies using surveys. Much of this research builds on advances in survey methodology, allowing researchers to uncover the perceptions, attitudes, beliefs and reasoning underlying policy views (See among others \cite{dechezlepretre2025fighting} for beliefs and reasoning about policies, \cite{haaland2023designing} for the design of information experiments and \cite{stantcheva2023run} for a review of the survey methodology). The survey methodology is applied in many settings, for example to understand people's perception of inflation  \parencite{binetti2024people} or to investigate the origins and implications of zero-sum thinking  \parencite{NBERw31688}). Most closely related \textcite{stantcheva2021understanding} explores beliefs, reasoning and support for tax policy.

Second, this project is related to the literature on tax enforcement and tax evasion, and especially to the literature on the relevance of information reporting. \textcite{slemrod2007cheating} has documented that misreported income in the US varies dramatically by income type. While income from wages, salaries, interest and dividends -- all of which are subject to third-party reporting -- is rarely misreported, income from self-employment -- typically self-reported -- has estimated tax gaps above 50\%. Similarly, \textcite{kleven2011unwilling} discover near-complete compliance and zero deterrence effects for third-party reported income sources in Denmark. And, \textcite{pomeranz2015no} looks at VAT paper trails in Chile and finds that the effects of randomized deterrence messages are much lower in areas covered by VAT paper trails, suggesting that evasion is much lower in areas covered by third-party reporting. Other prior work investigates effects of information mechanisms such as consumer receipt collection \parencite{naritomi2019consumers}, the role of digital payment reporting forms like 1099-K \parencite{slemrod2017does}, and discusses the theoretical foundations for why third-party reporting improves compliance \parencite{kleven2016can}. Finally, \textcite{jensen2022employment} documents how modern tax systems arise over time and finds evidence supporting the interpretation that the rise in third-party covered income (through increases in the employee share) drives expansions of the income tax base.

More recent research shows that the international Automatic Exchange of Information (AEoI) reduces offshore tax evasion \parencite{boas2024taxing,baselgia2023compliance}. Another key finding of the analysis of offshore tax evasion is that it is much more prevalent among the highest income earners. Contrary to the findings of \textcite{johns2010distribution}, which revealed that the net misreporting percentage generally increases with income group but peaks at the 99th percentile, \textcite{alstadsaeter2019tax} analyse leaked data and conclude that offshore tax evasion is highly concentrated among the super-rich. They estimate that the top 0.01\% of households evade approximately 25\% of their taxes. By contrast, tax evasion detected in stratified random tax audits is less than 5 percent throughout the distribution. These results are consistent with those of \textcite{boas2024taxing}, who analysed offshore tax evasion in Denmark, and \textcite{guyton2021tax}, who studied tax evasion at the top of the income distribution in the US. They showed that previous estimates had failed to capture sophisticated evasion via offshore accounts and pass-through businesses, therefore underestimating top 1\% tax evasion.

This project also relates to the literature on tax complexity, tax compliance costs and the government's role in providing services to reduce them. \textcite{craig2024tax} provide a unified analysis of taxation and taxpayer education when individuals have an incomplete understanding of the tax system. \textcite{feldman2016taxpayer} analyse whether tax complexity causes misperceptions about the marginal tax rate by looking at the effects of lump sum tax liability changes. \textcite{abeler2015complex} present evidence that people underreact to changes in tax incentives when facing more complex tax systems. While \textcite{zwick2021costs} explores how tax code complexity alters corporate behavior, \textcite{blumenthal1992compliance} and \textcite{slemrod1995income} analyse the compliance costs of complex tax system for the U.S. individual taxpayers. 

To understand how costly taxpayers in the US perceive filing declarations to be, \textcite{benzarti2020taxing} observes how taxpayers choose between itemizing deductions and claiming the standard deduction, finding that taxpayers forego large tax savings to avoid compliance costs. \textcite{hauck2024optional} shows that non-filing of tax declarations is particularly common at the lower end of the income distribution in Germany, which reduces effective redistribution. 

\textcite{blesse2019people} present survey evidence, eliciting the preferences of German taxpayers regarding the trade-off between simplifying the German tax code and reducing its progressivity. They discover that, while most taxpayers support simplifying the tax system, this support is partly driven by misconceptions about the trade-offs involved in tax complexity. \textcite{benzarti2024rising} ask US respondents about the complexity of the tax code and find that taxpayers perceive an increase in tax complexity and filing costs, and that the majority would be willing to pay for a simpler tax system with prepopulated tax returns.

Our survey is contributing to existing research in several ways. First, besides providing a more detailed measure for the support for prepopulated tax returns we also examine people’s demand for tax software. Second, by eliciting people’s views on high levels of third-party reporting and the international exchange of information, we provide novel evidence for the support of the use of tax enforcement tools involving information reporting. Third, we explore how this support interplays with support for government-provided tax filing services. For example, we investigate whether support for third-party reporting increases when taxpayers are aware that it is a prerequisite for government-provided prepopulated tax returns, or if third-party reported data were only used to reduce tax compliance costs. Finally, by eliciting important underlying considerations and beliefs, and by experimentally providing policy-relevant information, we investigate about which factual matters people might be misinformed and which considerations actually drive policy views. 

\section*{Notes}

First, this project relates to the literature on tax enforcement, tax evasion and especially the role of third-party reporting in reducing tax evasion. 
\begin{itemize}
    \item There is extensive research documenting the extent and distribution of tax evasion:  \textcite{johns2010distribution} shows that the net misreporting percentage generally increases with income group. It peaks at 99-99.5\%. 
    \item \textcite{slemrod2007cheating} shows that two thirds of all underreporting of income happens on the individual income tax. The understated income (80\%) is the major driver of the income tax gap. The most striking aspect of the aggregate tax gap estimates is the huge variation in the rate of misreporting as a percentage of true income by type of income. Almost none of the misreported income can be accounted to income sources such as wages, salaries, interest and dividends, which are all third party reported. In sharp contrast, self-employed income which is not subject to information reports, shows estimated tax gaps above 50\%. 
    \item \textcite{kleven2011unwilling} find that (in Denmark) tax evasion rate is close to zero for income subject to third-party reporting, but substantial for self-reported income. 
    \item \textcite{pomeranz2015no} looks at VAT-paper trails (in Chile) and finds that effects of randomized deterrence messages are much lower in areas that are covered in VAT paper trails (business to business), which suggests that evasion is much lower in areas covered under third party reporting (even if not directly linked to the government)
    \item \textcite{slemrod2017does} matches data from the 1099-K forms to what is reported by small business owners. They study the characteristics of firms by the ratio of receipts reported by Form 1099-K (K) to reported receipts(R)in 2011, the year of the policy change.
    \item \textcite{boas2024taxing} (unpublished) study the effects of AEoI on tax compliance by analyzing the universe of information reports sent by foreign banks to Danish authorities, matched to population-wide micro-data on income, wealth, and cross-border bank transfer. They estimate that automatic exchange of information has closed about 70\% of the offshore tax gap. 
    \item \textcite{baselgia2023compliance} (unpublished) looks at compliance effects of AEoI in Switzerland. Shows that tax evasion is more evenly distributed in Switzerland. 
    \item \textcite{alstadsaeter2019tax} show that offshore tax evasion is highly concentrated among the rich (In Scandinavia). The skewed distribution of offshore wealth implies high rates of tax evasion at the top: they find that the 0.01 percent richest households evade about 25 percent of their taxes. By contrast, tax evasion detected in stratified random tax audits is less than 5 percent throughout the distribution.
    \item \textcite{guyton2021tax} study tax evasion at the top of the income distribution and find that under the audit methods 2006-2013, individual random audit data failed to capture sophisticated evasion via offshore accounts and pass-through business.
    \item \textcite{kleven2016can} another study related to third party reporting: They develop an agency model explaining why third-party reporting by firms makes tax enforcement successful. 
    \item \textcite{naritomi2019consumers} another study looking at third party reporting: studies an anti-tax evasion program that rewards consumers for ensuring that firms report sales and establishes a verification system to aid whistle-blowing consumers in Sao Paulo, Brazil (Nota Fiscal Paulista).
    \item \textcite{jensen2022employment} another study looking at third party reporting: builds a new microdatabase that covers 100 countries to document how the modern tax system arises over development. Additional evidence supports the interpretation that the rise in third-party covered income through increases in employee share drives expansions of the income tax base over development.
    \item \textcite{slemrod2019tax} includes a review related to third party reporting 
\end{itemize}
\\
\begin{itemize}
    \item Tax simplicity and Complexity: 
    \begin{itemize}
        \item \textcite{craig2024tax} provide a unified analysis of taxation and taxpayer education when individuals have an incomplete understanding of the tax system.
        \item \textcite{feldman2016taxpayer} analyse whether tax complexity causes misperceptions about the marginal tax rate by looking at the effects of lump sum tax liability changes. 
        \item \textcite{abeler2015complex} analyse how people react to different levels of tax complexity and find that they underreact to changes in tax incentives when facing more complex tax system. 
        \item \textcite{blumenthal1992compliance} and \textcite{slemrod2005etiology} analyse the compliance cost of complex tax systems for the U.S. individual tax payers and \textcite{zwick2021costs} for U.S. corporate taxpayers. 
    \end{itemize}
    \item \textcite{goodman2023automated} argue that the IRS already receives information from third parties on many income sources reported on tax returns, including wages, unemployment compensation, interest, dividends, capital gains, nonemployee compensation, and income from partnerships and S corporations. Using two approaches they simulate prepopulated tax returns for 42 and 48\% of taxpayers successfully. For both approaches they find that success rates are decreasing in income. 
    \item \textcite{benzarti2020taxing} observes how US taxpayers choose between itemizing deductions and claiming standard deduction. Interpreting this as revealed preferences he finds that the cost of itemizing ranges from 0.6 percent to 0.8 percent of adjusted gross income (AGI); i.e., the disutility derived from itemizing is equivalent to working 10 to 15 hours, which is substantially larger than previous estimates.
    \item \textcite{benzarti2021estimating} relies on same approach from \textcite{benzarti2020taxing} to estimate tax filing costs but then also extrapolates the costs of filing to other schedules and forms over time. 
    \item \textcite{hauck2024optional} quantify the impact of optional non-filing on effective taxation of low-income taxpayers. They show that non-filers face over-withholding and pay more taxes than intended by the tax schedule. Non-filing increases effective tax rates at the bottom of the income distribution, because low-income taxpayers are most likely to refrain from optional filing while facing substantial tax over-withholdings. They show that non-filing is very common and particularly so at the lower end of the income distribution. Nonfilers over-remit taxes at all income levels and over-remittances are substantial with an annual value of at least 951 million. On the individual level, the average non-filer over-remits. And on average, non-filers’ effective ATRs are 1.2 percentage points higher than their statutory ATRs
    \item \textcite{blesse2019people} implemented a survey of German taxpayers eliciting their preferences over the trade-off between simplifying the German tax code and making it less progressive. Find that most of the taxpayers support simplification of the tax system, however that support is party driven by misinformation about trade-offs of tax complexity. 
\end{itemize}

Third, this project relates to the literature on measuring people's support for policies using experimental surveys. 
\begin{itemize}
    \item \textcite{stantcheva2021understanding} elicits preferences on tax policy in general. While she concentrates on the question which tax schedule of income and wealth taxes people prefer, we try to understand people's views on which instruments tax administration should use to enforce that tax schedule. 
    \item related but not very important: Literature on empirical determinants of the demand of redistribution. Review in Alesina and Giuliano, 2011 and also a nice short overview in Fehr, Epper and Senn 2023
    \item related but maybe not very important: Research related to tax compliance.  (See for example at literature of Ricardo Perez Truglia)
    \item \textcite{benzarti2024rising} elicit taxpayers perceptions about the complexity of the tax code. They find that the majority of their survey respondents believe that the tax code’s complexity makes the tax system less fair overall. And, the majority of respondents believe that increased complexity fosters evasion, further exacerbating the perception of unfairness that complexity imposes on taxpayers. According to their survey, the majority of taxpayers would be willing to pay for receiving a prepopulated tax return. And among those that are willing to pay for them, they would pay \$77 on average. They also asked US taxpayers how much time it takes them to file their taxes. Their estimates are lower than those of the IRS or Benzarti (2021). This could be because the sample of US taxpayers in our survey differs from that used by the IRS. And it is consistent with the fact that our respondents believe that it is less tedious to file taxes for themselves than for others.
\end{itemize}
 
Details: 
\begin{itemize}
    \item \textcite{benzarti2024rising} elicit taxpayers perceptions about the complexity of the tax code. They find that the majority of their survey respondents believe that the tax code’s complexity makes the tax system less fair overall. And, the majority of respondents believe that increased complexity fosters evasion, further exacerbating the perception of unfairness that complexity imposes on taxpayers. According to their survey, the majority of taxpayers would be willing to pay for receiving a prepopulated tax return. And among those that are willing to pay for them, they would pay \$77 on average. They also asked US taxpayers how much time it takes them to file their taxes. Their estimates are lower than those of the IRS or Benzarti (2021). This could be because the sample of US taxpayers in our survey differs from that used by the IRS. And it is consistent with the fact that our respondents believe that it is less tedious to file taxes for themselves than for others.
    \item \textcite{goodman2023automated} argue that the IRS already receives information from third parties on many income sources reported on tax returns, including wages, unemployment compensation, interest, dividends, capital gains, nonemployee compensation, and income from partnerships and S corporations. Using two approaches they simulate prepopulated tax returns for 42 and 48\% of taxpayers successfully. For both approaches they find that success rates are decreasing in income. 
    \item \textcite{benzarti2020taxing} observes how US taxpayers choose between itemizing deductions and claiming standard deduction. Interpreting this as revealed preferences he finds that the cost of itemizing ranges from 0.6 percent to 0.8 percent of adjusted gross income (AGI); i.e., the disutility derived from itemizing is equivalent to working 10 to 15 hours, which is substantially larger than previous estimates.
    \item \textcite{benzarti2021estimating} relies on same approach from \textcite{benzarti2020taxing} to estimate tax filing costs but then also extrapolates the costs of filing to other schedules and forms over time. 
    \item \textcite{blesse2019people} implemented a survey of German taxpayers eliciting their preferences over the trade-off between simplifying the German tax code and making it less progressive. Find that most of the taxpayers support simplification of the tax system, however that support is party driven by misinformation about trade-offs of tax complexity. 
    \item \textcite{hauck2024optional} quantify the impact of optional non-filing on effective taxation of low-income taxpayers. They show that non-filers face over-withholding and pay more taxes than intended by the tax schedule. Non-filing increases effective tax rates at the bottom of the income distribution, because low-income taxpayers are most likely to refrain from optional filing while facing substantial tax over-withholdings. They show that non-filing is very common and particularly so at the lower end of the income distribution. Nonfilers over-remit taxes at all income levels and over-remittances are substantial with an annual value of at least 951 million. On the individual level, the average non-filer over-remits. And on average, non-filers’ effective ATRs are 1.2 percentage points higher than their statutory ATRs
\end{itemize}




\newpage
\printbibliography
\end{document}